% Generated draft: requires mathematical and source review.
\documentclass[11pt]{article}
\usepackage[margin=25mm]{geometry}
\usepackage{fontspec}
\setmainfont{Latin Modern Roman}
\usepackage{amsmath,amssymb,mathtools,unicode-math}
\setmathfont{Latin Modern Math}
\usepackage{enumitem,needspace,xurl,hyperref}
\setlength{\parindent}{0pt}
\setlength{\parskip}{5pt}
\setlength{\emergencystretch}{3em}
\newcommand{\N}{\mathbb{N}}
\newcommand{\Z}{\mathbb{Z}}
\newcommand{\Q}{\mathbb{Q}}
\newcommand{\R}{\mathbb{R}}
\newcommand{\C}{\mathbb{C}}
\newcommand{\F}{\mathbb{F}}
\newcommand{\A}{\mathbb{A}}
\newcommand{\PP}{\mathbb{P}}
\newcommand{\E}{\mathbb{E}}
\newcommand{\eps}{\varepsilon}
\newcommand{\dd}{\,\mathrm{d}}
\begin{document}
\section*{UnsolvedMath ID: 801032}
\textbf{Source problem code:} catalogue entry
\subsection{Definitions and mathematical statement}
A Polish space is a separable topological space admitting a compatible complete metric. A metric presentation consists of a complete separable metric space \((X,d)\) together with a numbered dense sequence \((x_n)_{n\in\N}\). Such a presentation is left-c.e. if there is a uniformly computable family of rational numbers \(a_{i,j,s}\), nondecreasing in \(s\), such that
\[
\lim_{s\to\infty}a_{i,j,s}=d(x_i,x_j)\qquad(i,j\in\N).
\]
It is right-c.e. if the analogous rational approximations are nonincreasing. It is computable if a single algorithm, on input \(i,j,k\in\N\), produces a rational \(q\) satisfying \(\lvert q-d(x_i,x_j)\rvert\leq 2^{-k}\). A space is compact if every open cover has a finite subcover. The supplied source does not define its additional effectiveness requirement in the phrase ``effectively compact''; that requirement must be specified relative to the presentation before the effective version can be made precise.

The source asks: Does every \emph{(effectively) compact left-c.e. Polish space} have a \emph{computable (right-c.e.) Polish copy}? A quantified template for the question is
\[
\begin{gathered}
\text{For every left-c.e. presentation }(X,d,(x_n)),\\
\mathcal K(X,d,(x_n))\ \Longrightarrow\
\exists (Y,\rho,(y_n))\,
\bigl[\mathcal P(Y,\rho,(y_n))\ \land\ X\sim Y\bigr],
\end{gathered}
\]
where the existential quantifier ranges over complete separable metric presentations. Here \(\mathcal K\) denotes the intended compactness or effective compactness condition, and \(\mathcal P\) denotes the intended computable or right-c.e. presentation condition. The parenthetical wording does not specify how these alternatives are paired, so this template is not an assertion that all four combinations are intended.

There is also an unresolved interpretation of \(X\sim Y\). If ``copy'' means an isometric copy of the given metric space, the requirement is a bijection \(f:X\to Y\) with
\[
\rho(f(x),f(x'))=d(x,x')\qquad(x,x'\in X).
\]
If it means a topological copy, the requirement is a homeomorphism \(f:X\to Y\). These requirements are different, and the supplied record does not explicitly select one. Neither interpretation imposes computability on \(f\) unless an additional requirement is supplied.

As background, the record attributes to Melnikov and Ng the existence of a left-c.e. Polish space with no computable Polish presentation, and to Koh, Melnikov, and Ng the existence of a left-c.e. Polish space with no right-c.e. Polish presentation. It also attributes to Melnikov and Ng the positive result that every left-c.e. Stone space is homeomorphic to a computable Polish space. A Stone space is a compact Hausdorff space with a basis of clopen sets. These are supplied background claims; their precise presentation conventions and hypotheses are not independently verified here.
\subsection{Short English statement}
The question asks whether compactness, possibly with an effectiveness condition, ensures that a left-c.e. Polish space has a computable or right-c.e. Polish copy. The source leaves the pairing of these alternatives, the effective compactness convention, and whether a copy must preserve distances or only topology unresolved. The registry labels this problem open; this draft does not independently verify that status.
\subsection{Sources}
\begin{itemize}
\item Supplied local UnsolvedMath registry record, permanent ID \(801032\), exact source code \texttt{catalogue entry}, titled ``Open Problems in Mathematical Logic --- Problem 44.'' The record lists Barmpalias, Gao et al. as proposers, gives the year 2026, and labels the problem \texttt{open}.
\item Supplied source attribution: Barmpalias, Gao et al., \emph{Open Problems in Mathematical Logic} (2026), Problem 44, \url{https://arxiv.org/abs/2608.26628}. The background records access on 2026-10-16 and extraction by \texttt{arxiv\_latex}. This is supplied provenance, not a claim of independent access or verification.
\item Supplied background attributions: Melnikov and Ng for the nonexistence of a computable presentation in a left-c.e. example and for the positive Stone-space result; Koh, Melnikov, and Ng for the nonexistence of a right-c.e. presentation in a left-c.e. example. No separate bibliographic references for these results are supplied.
\end{itemize}
\end{document}
